% answers_short.tex —— 2005 年·全国卷Ⅰ·理科 · 简短答案速查
% 由 tools/make_short_answers.py 自 answers.tex 的 \gksolhead/\gkanswer 提取，
% 并与 2005/2005_数学真题_全国1理_答案要点.md 逐题比对后生成，勿手改（改源 MD 后重跑）。

\begin{gknotice}
\textbf{使用说明}：本节只给最终答案便于快速核对，完整推导见解析版。答案为编写者独立推导，\textbf{非官方标准答案}。
\end{gknotice}

\gkpart{选择题}
\begin{center}\begin{tabular}{ccccccccccccc}
\toprule
题号 & 1 & 2 & 3 & 4 & 5 & 6 & 7 & 8 & 9 & 10 & 11 & 12 \\
\midrule
答案 & A & C & B & C & A & D & C & C & C & B & B & D \\
\bottomrule
\end{tabular}\end{center}

\gkpart{填空题}
\begin{itemize}[leftmargin=2.6em]
  \item[\textbf{13.}] 155
  \item[\textbf{14.}] 672
  \item[\textbf{15.}] 1
  \item[\textbf{16.}] ①②④
\end{itemize}

\gkpart{解答题}
\begin{itemize}[leftmargin=2.6em]
  \item[\textbf{17.}] (1) $\varphi=-\dfrac{3\pi}{4}$；(2) $x\in\left(\dfrac{\pi}{8}+k\pi,\dfrac{5\pi}{8}+k\pi\right)\ (k\in\mathbb Z)$；(3) 不存在相切情形
  \item[\textbf{18.}] (1) 面 $PAD\perp$ 面 $PCD$；(2) $\cos\theta=\dfrac{2}{\sqrt{10}}=\sqrt{\dfrac25}$；(3) 两平面的锐角夹角满足 $\cos\theta=\dfrac23$
  \item[\textbf{19.}] (1) $q>-1$；(2) $T_n=q^2\left(q-\dfrac32\right)S_n$。令 $\alpha$ 为方程 $q^2(q-\dfrac32)=1$ 在 $(-1,+\infty)$ 内的唯一根，则 $S_n>T_n$（$-1<q<\alpha$），$S_n=T_n$（$q=\alpha$），$S_n<T_n$（$q>\alpha$）
  \item[\textbf{20.}] (1) $P(\xi=0)=\dfrac{343}{512},\ P(\xi=10)=\dfrac{147}{512},\ P(\xi=20)=\dfrac{21}{512},\ P(\xi=30)=\dfrac1{512}$；(2) $E(\xi)=3.75$
  \item[\textbf{21.}] (1) $e=\dfrac{\sqrt6}{3}$；(2) $\lambda^2+\mu^2=1$
  \item[\textbf{22.}] (1) 最小值为 $-1$，在 $x=\dfrac12$ 时取得；(2) 不等式成立，且等号当且仅当 $p_1=\cdots=p_{2n}=\dfrac1{2n}$
\end{itemize}
